% =====================================================================
%  CSCI 3130 — Formal Languages & Automata Theory
%  Full Midterm Cheat Sheet (Lectures 1-6, Ch. 1-6, CYK)
%  Course: 何宗易 Tsung-Yi Ho (CUHK)
%  Sources: github.com/gingersea/3130 — lec1..lec6.pdf, quiz1.pdf, Print.pdf
%  Layout: A4, 9pt, three columns, landscape, 4 pages (double-sided A4 x2)
% =====================================================================
\documentclass[9pt,a4paper,landscape]{extarticle}
\usepackage[margin=0.5cm,top=0.6cm,bottom=0.6cm]{geometry}
\usepackage{multicol}
\usepackage{enumitem}
\usepackage{amsmath,amssymb}
\usepackage[T1]{fontenc}
\usepackage{lmodern}
\usepackage{microtype}
\usepackage[svgnames]{xcolor}
\usepackage{tcolorbox}
\usepackage{fancyhdr}
\usepackage{titlesec}

\setlist{nosep,leftmargin=1.1em,topsep=1pt,itemsep=0pt}
\setlength{\columnseprule}{0.4pt}
\setlength{\columnsep}{8pt}
\setlength{\parindent}{0pt}
\setlength{\parskip}{0pt}
\raggedcolumns

\pagestyle{fancy}
\fancyhf{}
\fancyhead[L]{\scriptsize\color{SteelBlue}\textbf{CSCI 3130 -- Formal Languages \& Automata Theory -- Full Midterm Cheat Sheet}}
\fancyhead[R]{\scriptsize\textbf{Page \thepage/4}}
\fancyfoot[C]{\scriptsize Sources: lec1--lec6 (Tsung-Yi Ho, CUHK) \textbullet\ github.com/gingersea/3130}

\titleformat{\section}{\color{SteelBlue}\normalsize\bfseries}{}{0pt}{}[\vspace{-2pt}{\color{SteelBlue}\titlerule[0.5pt]}]
\titlespacing{\section}{0pt}{4pt}{2pt}
\titlespacing{\subsection}{0pt}{3pt}{1pt}
\titleformat{\subsection}{\small\bfseries}{}{0pt}{}

\tcbset{colback=AliceBlue!45,colframe=SteelBlue!60!black,boxrule=0.5pt,
  arc=1.5mm,left=2.5pt,right=2.5pt,top=1.5pt,bottom=1.5pt,before skip=2pt,after skip=2pt}
\newtcbox{\frm}[1][SteelBlue]{on line,colback=#1!8,colframe=#1!60!black,size=small,
  arc=1pt,left=1.5pt,right=1.5pt,top=0.5pt,bottom=0.5pt,boxrule=0.4pt}
\newcommand{\T}[1]{\text{\scriptsize #1}}

\newcommand{\ex}[1]{\frm[FireBrick]{\textbf{Ex.\ #1}}}
\newcommand{\df}[1]{\frm[DodgerBlue]{\textbf{Def.\ #1}}}
\newcommand{\thm}[1]{\frm[ForestGreen]{\textbf{Thm.\ #1}}}

\begin{document}
\begin{multicols*}{3}

%======================================================================
\section{Ch 1 -- Mathematical Preliminaries (Lec 1)}
%----------------------------------------------------------------------
\subsection{Sets}
\df{set} collection of elements. $1\!\in\!A$, ship $\notin B$. Representations: explicit, ellipsis, set-builder $S=\{j: j>0,\ j=2k,\ k>0\}$.
Operations: $A\cup B$, $A\cap B$, $A-B$ (e.g.\ $\{1,2,3\}-\{2,3,4,5\}=\{1\}$), complement $\bar{A}=U-A$. $\bar{\bar{A}}=A$.
\textbf{DeMorgan:} $\overline{A\cup B}=\bar{A}\cap\bar{B}$; $\overline{A\cap B}=\bar{A}\cup\bar{B}$.
$\emptyset=\{\}$; $S\cup\emptyset=S$, $S\cap\emptyset=\emptyset$, $S-\emptyset=S$.
\textbf{KEY:} $\emptyset\neq\{\lambda\}$: $|\{\}|=0$, $|\{\lambda\}|=1$, $|\lambda|=0$.
Subset $A\subseteq B$ (proper $\subset$). Disjoint: $A\cap B=\emptyset$.
Cardinality $|A|$. Powerset $2^S=\mathcal P(S)=\{$all subsets$\}$; if $S$ finite then $|2^S|=2^{|S|}$.
\ex{1.1} $S=\{a,b,c\}$: $2^S=\{\emptyset,\{a\},\{b\},\{c\},\{a,b\},\{a,c\},\{b,c\},\{a,b,c\}\}$, $|2^S|=2^3=8$.
Cartesian product $A\times B=\{(x,y):x\in A,\ y\in B\}$, $|A\times B|=|A||B|$. Order matters: $(4,2)\in A\times B$ but $(2,4)\notin A\times B$.
Partition of $S$: $S_1,\ldots,S_n$ mutually disjoint, union $=S$, all nonempty.

%----------------------------------------------------------------------
\subsection{Functions}
$f:A\to B$ assigns each domain element a \emph{unique} range element. Total if defined on whole $A$; else partial.

%----------------------------------------------------------------------
\subsection{Asymptotic Analysis}
$f=O(g)$: $\exists c,n_0>0$: $0\le f(n)\le c\,g(n)\ \forall n\ge n_0$ (upper).
$f=\Omega(g)$: $0\le c\,g(n)\le f(n)$ (lower).
$f=\Theta(g)$ iff both $O$ and $\Omega$ (tight).
\ex{1.3} $f=2n^2+3n,\ g=n^3,\ h=10n^2+100$: $f=O(g)$, $g=\Omega(h)$, $f=\Theta(h)$.

%----------------------------------------------------------------------
\subsection{Relations}
Relation: a set of pairs, $x_i\,R\,y_i$; one $x$ may relate to many $y$ (unlike functions). Equivalence relation: \textbf{reflexive} ($xRx$), \textbf{symmetric} ($xRy\Rightarrow yRx$), \textbf{transitive} ($xRy\wedge yRz\Rightarrow xRz$).
\ex{1.4} $x\equiv y$ iff $x\bmod 3=y\bmod 3$ on $\mathbb N$: $2\equiv 5$, $12\equiv 0$, $0\equiv 36$; an equivalence (equivalence classes: residues mod 3).

%----------------------------------------------------------------------
\subsection{Graphs}
Digraph $G=(V,E)$: nodes $V$, edges $E\subseteq V\times V$ (ordered pairs).
\df{walk} sequence of adjacent edges.
\df{path} walk with no edge repeated; \emph{simple path}: no node repeated.
\df{cycle} walk from a node back to itself; \emph{simple cycle}: only base node repeated. \df{loop} edge from vertex to itself.
\textbf{Euler tour:} cycle containing each \emph{edge} once. \textbf{Hamiltonian cycle:} simple cycle containing all \emph{nodes}.
Finding all simple paths: grow from origin, depth-first, never repeat a node.
\df{tree} graph with no cycles. root, parent/child, leaf, level (root = level 0), height = longest root--leaf path.

%----------------------------------------------------------------------
\subsection{Proof Techniques}
\begin{enumerate}
\item \textbf{Direct/constructive:} assume $X$, show $Y$.
\item \textbf{Induction:} basis $P_1,\dots,P_b$ true; hypothesis $P_1,\dots,P_k$; step: show $P_{k+1}$.
\ex{1.5} Binary tree of height $n$ has $\le 2^n$ leaves. Let $L(i)$ = max leaves at height $i$. Basis $L(0)=1$. Hyp. $L(i)\le 2^i$, $i\le k$. Step: each level-$k$ leaf adds at most 2 children: $L(k+1)\le 2L(k)\le 2\cdot 2^k=2^{k+1}$. (Note: recursion, e.g.\ $f(n)=f(n{-}1)+f(n{-}2)$, is a different thing.)
\ex{1.6} $n$ mutually intersecting lines divide plane into $A(n)$ regions; $A(1)=2,A(2)=4,A(3)=7$, $A(n+1)=A(n)+n+1$; solved form $A(n)=\tfrac{n(n+1)}{2}+1$.
\item \textbf{Contradiction:} assume $P$ false, derive impossible conclusion.
\ex{1.7} $\sqrt2\notin\mathbb Q$: assume $\sqrt2=n/m$ coprime; $2m^2=n^2$ so $n$ even, $n=2k$; $2m^2=4k^2\Rightarrow m^2=2k^2$ so $m$ even; $n,m$ share factor 2 -- contradiction.
\end{enumerate}

%----------------------------------------------------------------------
\subsection{Languages \& Grammars (three basic concepts: languages, grammars, automata)}
Alphabet $\Sigma$ (finite symbols); string $w=a_1\cdots a_n$ over $\Sigma$; language = any subset of $\Sigma^*$.
\begin{itemize}
\item Concatenation $wv$; Reverse $w^R=a_n\cdots a_1$; Length $|w|$, $|uv|=|u|+|v|$.
\item Empty string $\lambda$: $|\lambda|=0$, $\lambda w=w\lambda=w$.
\item Substring = consecutive chars; prefix/suffix: $w=uv$ ($u$ prefix, $v$ suffix).
\item Power $w^n=ww\cdots w$ ($n$ times), $w^0=\lambda$; $(abba)^2=abbaabba$.
\item $\Sigma^*$ = all strings over $\Sigma$ (incl.\ $\lambda$); $\Sigma^+=\Sigma^*-\{\lambda\}$.
\end{itemize}
\textbf{Language operations:} union, intersection, difference, complement $\bar L=\Sigma^*-L$, reverse $L^R=\{w^R:w\in L\}$, concatenation $L_1L_2=\{xy:x\in L_1,y\in L_2\}$, power $L^n$ ($L^0=\{\lambda\}$), \textbf{Kleene star} $L^*=L^0\cup L^1\cup L^2\cup\cdots$, positive closure $L^+=L^*-\{\lambda\}$.
\ex{} $\{a,ab,ba\}\{b,aa\}=\{ab,aaa,abb,abaa,bab,baaa\}$.
\ex{} $L=\{a^nb^n:n\ge0\}$, then $L^2=\{a^nb^na^mb^m:n,m\ge0\}$ -- $n,m$ \emph{unrelated} ($aabbaaabbb\in L^2$).

\df{1.1 grammar} 4-tuple $G=(V,T,S,P)$: $V$ variables, $T$ terminals, $S\in V$ start variable, $P$ productions $x\to y$ with $x\in(V\cup T)^+$, $y\in(V\cup T)^*$.
\textbf{Derivation:} $\alpha\Rightarrow\gamma$ in one step iff $\alpha=\alpha_1A\alpha_2$, $A\to\beta\in P$, $\gamma=\alpha_1\beta\alpha_2$. $\alpha\overset{*}{\Rightarrow}\gamma$: zero or more steps.
\df{1.2} $L(G)=\{w\in T^*: S\overset{*}{\Rightarrow} w\}$; $S,w_1,\dots,w_n$ are \emph{sentential forms}; $w$ is a \emph{sentence}.
\ex{1.11} $S\to aSb\ |\ \lambda$: $S\Rightarrow aSb\Rightarrow aaSbb\Rightarrow aabb$; $L(G)=\{a^nb^n:n\ge0\}$. $aabb$ sentence, $aaSbb$ sentential form.
\ex{1.12} Grammar for $L=\{a^nb^{n+1}:n\ge0\}$: $S\to Ab,\ A\to aAb,\ A\to\lambda$ (one extra $b$).
\ex{1.13} $S\to SS|\lambda|aSb|bSa$ generates $\{w: n_a(w)=n_b(w)\}$; proof by induction on $|w|=2n$: $w=aw_1b$ (or $bw_1a$): $S\Rightarrow aSb\overset{*}{\Rightarrow}aw_1b$; else $w=w_1w_2$, $S\Rightarrow SS\overset{*}{\Rightarrow}w_1w_2$.
\textbf{Equivalent grammars:} $L(G_1)=L(G_2)$.
\ex{1.14} $G_1$: $S\to aSb|\lambda$; $G_2$: $S\to aAb|\lambda,\ A\to aAb|\lambda$ are equivalent (both $\{a^nb^n\}$).

%======================================================================
\section{Ch 2 -- Finite Automata (Lec 2)}
%----------------------------------------------------------------------
\subsection{DFA}
\df{2.1 DFA} $M=(Q,\Sigma,\delta,q_0,F)$: finite $Q$; input alphabet $\Sigma$; transition function $\delta:Q\times\Sigma\to Q$ (\textbf{total}, exactly one move per symbol); initial state $q_0\in Q$; final states $F\subseteq Q$.
\df{2.2 (extended $\delta^*$)} $\delta^*:Q\times\Sigma^*\to Q$ recursively:
$\delta^*(q,\lambda)=q$, $\ \delta^*(q,wa)=\delta(\delta^*(q,w),a)$.
\thm{2.1} There is a walk from $q_i$ to $q_j$ labeled $w$ iff $\delta^*(q_i,w)=q_j$.
\df{} $L(M)=\{w\in\Sigma^*:\delta^*(q_0,w)\in F\}$ (regular language).
$\bar{L}(M)=\{w:\delta^*(q_0,w)\notin F\}$.
\textbf{Trap state:} non-final state with self-loops on all symbols (absorbs bad prefixes).
\ex{2.2} $L=\{a^nb:n\ge0\}$: $q_0\overset{a}{\circlearrowleft}$, $q_0\overset{b}{\to}q_1$ (final), $q_1\overset{a,b}{\to}q_2$ trap.
\ex{2.3} prefix $ab$: $q_0\overset{a}{\to}q_1\overset{b}{\to}q_2$ (final loop); trap $q_3$.
\ex{2.4} without substring 001: states track longest 001-prefix seen: $\lambda,0,00,001$ (trap).
\ex{2.5} $L=\{awa:w\in\{a,b\}^*\}$ is regular (finite memory: remember first letter $a$; reject if other letters differ).
\ex{2.6} $L^2=\{aw_1aaw_2a\}$ likewise regular.
\textbf{Non-regular} (preview): $\{a^nb^n:n\ge0\}$ -- needs unbounded counting; proved in Ch 4.
\textbf{Regular language:} $L$ is regular iff some DFA $M$ has $L=L(M)$.
\textbf{Design drills} (states = memory of what matters): exactly four 0s (counter 0--4+trap); exactly four 0s and $\ge2$ 1s (5$\times$3 grid); length $\le 5$; ends in 1101; binary value divisible by 5 (state = remainder mod 5); contains 0101; avoids 110.
\ex{Ex 2.1.7(a)} $|w|\bmod 3=0$: 3-state cycle on $a,b$, final state 0.
\ex{Ex 2.1.9(a)} every 00 followed immediately by 1: states $\lambda,0,00$; from 00 only $0$ allowed; 00 then 1 returns.
\textbf{Other drills:} starts with 0 and odd length, or starts with 1 and even length; every odd position is 1.

%----------------------------------------------------------------------
\subsection{NFA}
\df{} $M=(Q,\Sigma,\delta,q_0,F)$ with $\delta:Q\times(\Sigma\cup\{\lambda\})\to 2^Q$. Differences from DFA: range is powerset $2^Q$; $\lambda$-transitions allowed (read head does not move); $\delta(q_i,a)$ may be $\emptyset$ (automaton \emph{hangs}).
\textbf{Accept:} string $w$ accepted iff \emph{some} computation consumes all of $w$ and ends in a final state. \textbf{Reject:} \emph{every} computation either ends non-final after consuming input, or hangs before input consumed.
$\lambda$ never appears on the input tape.
\df{2.5} $\delta^*(q_i,w)$ = set of states reachable via a walk labeled $w$ ($\lambda$-edges cost nothing).
\df{2.6} $L(M)=\{w\in\Sigma^*:\delta^*(q_0,w)\cap F\neq\emptyset\}$.
\ex{2.8} $M$: $q_0\overset{1}{\to}q_1$, $q_1\overset{0}{\to}\{q_0,q_2\}$, $q_0\overset{\lambda}{\to}\{q_0,q_2\}$: $L(M)=\{10\}^*$; $\delta(q_2,1)=\emptyset$.
\ex{2.9} walks: $\delta^*(q_1,a)=\{q_0,q_1,q_2\}$; walk length may exceed $|w|$ (through $\lambda$-edges).
\ex{} $\{a^3\}\cup\{a^{2n}:n\ge1\}$ is easy for an NFA (two branches), tedious for a DFA -- nondeterminism gives concise descriptions, e.g.\ $S\to aSb|\lambda$.
\textbf{Exercise 2.3.7:} any NFA $\to$ equivalent NFA with a \emph{single} final state (add new final state, $\lambda$-edges from all old finals). Extreme case: NFA with no final state accepts $\emptyset$; add an unreachable final state.

%----------------------------------------------------------------------
\subsection{Equivalence of DFA and NFA}
\df{2.7} $M_1,M_2$ equivalent iff $L(M_1)=L(M_2)$.
\thm{2.2} Languages accepted by NFAs = regular languages (DFAs). Two directions:
(1) DFA $\subseteq$ NFA: trivial (DFA is a special NFA).
(2) NFA $\to$ DFA: \textbf{subset construction}:
\begin{enumerate}
\item Initial DFA state: $\{q_0\}$.
\item For DFA state $\{q_i,q_j,\dots,q_m\}$ and symbol $a$: compute union of all $\delta^*(q_\ell,a)$ in NFA; add DFA transition $\delta'(\{q_i,q_j,\dots\},a)=\{q'_i,q'_j,\dots\}$.
\item Repeat step 2 for all letters until no new states appear.
\item DFA state is final iff it \emph{contains} an NFA final state.
\end{enumerate}
DFA states $\in 2^Q$: $\emptyset,\{q_0\},\{q_1\},\{q_1,q_2\},\dots$ (possibly exponentially many; at most $2^{|Q|}$).
\textbf{Correctness proof (Thm 2.2):} induction on $|v|$: if in NFA $q_0\overset{v}{\Rightarrow}q_m$, then in DFA $\{q_0\}\overset{v}{\Rightarrow}\{q_m,\ldots\}\ni q_m$. Then $w\in L(M)\Rightarrow w\in L(M')$; converse similar; hence $L(M)=L(M')$.
\ex{2.12} NFA $q_0\overset{a}{\to}q_1$, $q_1\overset{\lambda}{\to}q_2$, $q_2\overset{b}{\circlearrowleft}$, $q_1$ final. DFA: $\{q_0\}\overset{a}{\to}\{q_1,q_2\}$ (final, since $q_1\in F$); $\{q_0\}\overset{b}{\to}\emptyset$; $\{q_1,q_2\}\overset{a}{\to}\emptyset$, $\overset{b}{\to}\{q_2\}$; $\{q_2\}\overset{b}{\to}\{q_2\}$; $\emptyset$ loops on $a,b$. $L(M')=\{ab^*\}=L(M)$.

%----------------------------------------------------------------------
\subsection{State Reduction (minimization)}
\df{2.8} States $p,q$ of a DFA are \textbf{indistinguishable} if for \emph{all} $w\in\Sigma^*$: $\delta^*(p,w)\in F\iff\delta^*(q,w)\in F$. Otherwise \textbf{distinguishable by} $w$ (some $w$ sends one to $F$, the other out).
Algorithm (partition refinement):
\begin{enumerate}
\item Remove \emph{inaccessible} states (unreachable from $q_0$).
\item Start: partition into $F$ and $Q-F$ (final vs non-final). Every pair across this cut is distinguishable by $w=\lambda$.
\item Split a block: if two states of the same block go (on the same symbol) to different blocks, they are distinguishable -- split accordingly. Repeat until no split possible.
\item Blocks = states of the minimal DFA (merge each block into one state).
\end{enumerate}
\ex{2.15} $F=\{q_0,q_1,q_3\}$, $\{q_2,q_4\}$ non-final. $\delta(q_0,0)=q_1$ vs $\delta(q_1,0)=q_2$ ($q_2\notin F$): $q_0,q_1$ distinguishable by 0; split $\{q_0,q_1,q_3\}\to\{q_0\},\{q_1,q_3\}$. Similarly $\delta(q_2,0)=q_1$ vs $\delta(q_4,0)=q_4$: split $\{q_2,q_4\}\to\{q_2\},\{q_4\}$.
Minimal DFA is unique (up to state renaming) -- gives algorithm for checking DFA equivalence: minimize both, compare.

%======================================================================
\section{Ch 3 -- Regular Expressions \& Regular Grammars (Lec 3)}
%----------------------------------------------------------------------
\subsection{Regular Expressions}
An RE is a \emph{notation}, not a language: it \emph{describes} a language (say ``$L$ is described by $(a+b+c)^*$'', never ``$L=(a+b+c)^*$'').
\df{3.1} Primitive: $\emptyset,\ \lambda$, and $a\in\Sigma$. If $r_1,r_2$ are REs, so are $r_1+r_2,\ r_1\cdot r_2,\ r_1^*,\ (r_1)$. Nothing else.
Semantics: $L(\emptyset)=\emptyset$; $L(\lambda)=\{\lambda\}$; $L(a)=\{a\}$; $L(r_1+r_2)=L(r_1)\cup L(r_2)$; $L(r_1r_2)=L(r_1)L(r_2)$; $L(r^*)=(L(r))^*$; $L((r))=L(r)$.
\ex{3.1} $(a+b\cdot c)^*\cdot(c+\emptyset)$ is an RE; $(a+b+)$ is not ($a^+=aa^*$).
\ex{3.2} $(a+b)\cdot a^*$ describes $\{a,aa,\dots,b,ba,\dots\}$.
\textbf{Operator priority:} $*>\cdot>+$. $a\cdot b+c$ means $(a\cdot b)+c$: $\{ab,c\}\neq\{ab,ac\}$.
\ex{3.3} $(a+b)^*(a+bb)$ = strings on $\{a,b\}$ ending in $a$ or $bb$.
\ex{3.4} $(aa)^*(bb)^*b=\{a^{2n}b^{2m+1}:n,m\ge0\}$.
\ex{3.5} contains 00: $(0+1)^*00(0+1)^*$.
\ex{3.6} no consecutive 0s: $r_1=(1+01)^*(0+\lambda)$; also $r_2=(1^*011^*)^*(0+\lambda)+1^*(0+\lambda)$. Infinitely many REs per language. Equivalent REs: $L(r_1)=L(r_2)$.
\textbf{More:} $L_2=\{x:|x|\ \text{even}\}$: $((0+1)(0+1))^*$; $L_3=\{x:x$ does not end in 01$\}$: $\lambda+0+1+(0+1)^*(00+10+11)$; $L_4$ = odd number of 0s: $1^*01^*(01^*01^*)^*$; exactly one $a$ on $\{a,b,c\}$: $(b+c)^*a(b+c)^*$.

%----------------------------------------------------------------------
\subsection{RE $\Leftrightarrow$ Regular Languages (Kleene's theorem)}
\thm{3.1} ($\Rightarrow$) $L(r)$ regular for every RE $r$ -- induction on structure: primitives $\emptyset,\lambda,a$ have trivial NFAs; union/concatenation/star use the standard NFA constructions (with single final states).
\thm{3.2} ($\Leftarrow$) every regular language has an RE -- from DFA/NFA via \textbf{generalized transition graphs (GTG)}: edges labeled by REs; merge parallel edges with $+$; label missing edges $\emptyset$.
\textbf{GTG state elimination} (NFA $\to$ RE procedure):
\begin{enumerate}
\item Single final state (else add one, $\lambda$ from old finals).
\item Complete the GTG (add $\emptyset$ edges).
\item If only $q_0$ and final $q_f$: $r=r_{ii}^*r_{ij}(r_{jj}+r_{ji}r_{ii}^*r_{ij})^*$.
\item Else pick $q_k$ to remove; for each pair $p,q\in\{i,j\}$: replace $r_{pq}\mapsto r_{pq}+r_{pk}r_{kk}^*r_{kq}$; delete $q_k$.
\item Simplify with $r+\emptyset=r,\ r\emptyset=\emptyset,\ \emptyset^*=\lambda$.
\item Repeat until two states.
\end{enumerate}
\ex{3.8} GTG $q_0\overset{a}{\circlearrowleft}\overset{b}{\to}q_1$, $q_1\overset{a+b}{\circlearrowleft}\overset{b}{\to}q_2$: eliminate $q_1$: $q_0\to q_2$: $bb^*(a+b)$; $q_0\to q_0$: $bb^*a$; result $r=(bb^*a)^*bb^*(a+b)b^*$.
\ex{3.11} $L=\{w:n_a(w)$ even, $n_b(w)$ odd$\}$: build 4-state NFA (EE,EO,OE,OO; start EE, final OE); eliminate to get RE (slides: from EE to OE via $\lambda$: with self-loops $aa$ and $bb$; answer $(aa+bb)^*\,(ab+ba)\,(aa+bb)^*\cdot(\ldots)$ -- compute via elimination; practice!).

%----------------------------------------------------------------------
\subsection{Regular Grammars}
\df{linear grammar} $\le1$ variable on RHS of each production. ($S\to aSb$, $S\to\lambda$ linear.)
\df{right-linear} $A\to xB$ or $A\to x$; \df{left-linear} $A\to Bx$ or $A\to x$ ($x\in T^*$). \textbf{Regular grammar} = right- or left-linear.
\textbf{Every regular grammar is linear; not conversely.} $G_3$: $S\to A,\ A\to aB|\lambda,\ B\to Ab$ is linear but not regular; $L(G_3)=\{a^nb^n\}$.
\ex{3.13} $G_1$: $S\to abS|a$ gives $L=(ab)^*a$. $G_2$: $S\to Aab,\ A\to Aab|B,\ B\to a$ gives $L=aab(ab)^*$.
\thm{3.3,3.4} Regular grammars $\Leftrightarrow$ regular languages.
\textbf{Right-linear grammar $\to$ NFA:} states = variables $V_0,V_1,\dots$ + special final state $V_F$ (start $V_0$). $V_i\to a_1a_2\cdots a_mV_j$: chain of $m$ edges via new intermediate nodes to $V_j$. $V_i\to a_1\cdots a_m$: chain to $V_F$.
\ex{3.15} $V_0\to aV_1$, $V_1\to abV_0|b$: NFA $V_0\overset{a}{\to}V_1$, $V_1\overset{a}{\to}x\overset{b}{\to}V_0$, $V_1\overset{b}{\to}V_F$; $L=(aab)^*ab$.
\textbf{Left-linear grammar $\to$ right-linear:} reverse RHS strings: $A\to Ba_1\cdots a_k\mapsto A\to a_k\cdots a_1B$ ($=A\to v^RB$); $A\to a_1\cdots a_k\mapsto A\to a_k\cdots a_1$. Then $L(G')=L(G)^R$; since $L(G')$ regular, $L(G)$ regular.
\textbf{NFA $\to$ right-linear grammar (BFS):} transition $q\overset{a}{\to}p$ gives production $q\to ap$; final state $q_f$ gives $q_f\to\lambda$. Then $L(G)=L(M)$.
\ex{} NFA $q_0\overset{a}{\to}q_1,\ q_1\overset{b}{\circlearrowleft},\ q_1\overset{a}{\to}q_2,\ q_2\overset{b}{\to}q_3,\ q_3\overset{\lambda}{\to}q_1,\ q_3$ final: productions $q_0\to aq_1,\ q_1\to bq_1|aq_2,\ q_2\to bq_3,\ q_3\to q_1|\lambda$; $L=ab^*ab(b^*ab)^*$.
\textbf{Summary:} four equivalent formalisms: DFA, NFA, RE, RG.

%======================================================================
\section{Ch 4 -- Properties of Regular Languages (Lec 4)}
%----------------------------------------------------------------------
\subsection{Closure Properties}
Regular languages are closed under: union, concatenation, star, complement, intersection, difference, reversal. Constructions assume NFAs with single final state.
\begin{itemize}
\item \textbf{Union} $L_1\cup L_2$: new start with $\lambda$-edges into both $M_1,M_2$; new single final.
\item \textbf{Concatenation} $L_1L_2$: $\lambda$ from $M_1$ final to $M_2$ start; final states of $M_2$.
\item \textbf{Star} $L_1^*$: new start (also final, to include $\lambda$) $\overset{\lambda}{\to}M_1$, $M_1$ final $\overset{\lambda}{\to}$ back to $M_1$ start; plus $\lambda$ to new final.
\item \textbf{Reversal} $L^R$: reverse all transitions; initial state becomes final; old final(s) become the single new initial (add new start with $\lambda$ to all old finals).
\item \textbf{Complement} $\bar L$: requires a \textbf{DFA} (complete it!); swap $F\leftrightarrow Q-F$. ($L=\{a^nb:n\ge0\}$ complemented: trap becomes final, etc.)
\item \textbf{Intersection} $L_1\cap L_2$: \textbf{product construction} $M=(Q_1\times Q_2,\dots)$: $\delta((p,q),a)=(\delta_1(p,a),\delta_2(q,a))$; $F=\{(p,q):p\in F_1,\ q\in F_2\}$. (Also: DeMorgan $\overline{\bar L_1\cup\bar L_2}$.)
\item \textbf{Difference} $L_1-L_2$: same product machine with $F=\{(p,q):p\in F_1,\ q\notin F_2\}$ (equivalently $L_1\cap\bar L_2$).
\end{itemize}
\ex{} $L_1=\{x: 00$ not substring$\}$, $L_2=\{x: x$ ends in 01$\}$: product DFA on states $AX,\dots,CZ$; final for $\cap$: only $F_1\times F_2$ pairs.
\df{4.1 homomorphism} $h:\Sigma\to\Gamma^*$; extend $h(w)=h(a_1)\cdots h(a_n)$, $h(L)=\{h(w):w\in L\}$.
\ex{4.2} $h(a)=ab,\ h(b)=bbc$: $h(aba)=abbbcab$, $h(\{aa,aba\})=\{abab,abbbcab\}$.
\ex{4.3} $L$ denoted $r=(a+b^*)(aa)^*$, $h(a)=dbcc,h(b)=bdc$: image denoted $(dbcc+(bdc)^*)(dbccdbcc)^*$.
\thm{4.3} If $L$ regular, $h(L)$ regular (substitute in the RE/transitions).

%----------------------------------------------------------------------
\subsection{Elementary Decision Questions}
Standard representations: DFA, NFA, RE, RG. Given DFA for $L$:
\begin{itemize}
\item \textbf{Membership} ($w\in L$?): simulate DFA on $w$; check final state.
\item \textbf{Emptiness} ($L=\emptyset$?): any path from $q_0$ to a final state (graph reachability)? No $\Rightarrow$ empty.
\item \textbf{Finiteness} ($L$ finite?): $L$ infinite iff there is a walk from $q_0$ to a final state containing a \emph{cycle}. (No cycle on any accepting path $\Rightarrow$ finite.)
\item \textbf{Equality} ($L_1=L_2$?): check $(L_1\cap\bar L_2)\cup(\bar L_1\cap L_2)=\emptyset$ (i.e.\ both differences empty). Equivalent: minimize both DFAs, compare.
\end{itemize}

%----------------------------------------------------------------------
\subsection{Pumping Lemma (identifying nonregular languages)}
\textbf{Pigeonhole principle:} $n$ pigeons, $m<n$ holes $\Rightarrow$ some hole has $\ge 2$ pigeons. In a DFA with $m$ states, any walk labeled by $|w|\ge m$ visits $>m$ states $\Rightarrow$ a state repeats. First repetition occurs within first $m$ moves.
\thm{4.8 (Pumping Lemma)} If $L$ is an infinite regular language, $\exists$ integer $m$ such that every $w\in L$, $|w|\ge m$, can be written $w=xyz$ with $|xy|\le m$, $|y|\ge1$, and $xy^iz\in L$ for all $i=0,1,2,\ldots$
\textbf{The game (adversary = Opponent $O$, you = Prover $P$):}
\begin{enumerate}
\item $O$ picks $m$ (number of states).
\item $P$ picks $w\in L$, $|w|\ge m$ (any $w$, but it must be \emph{in} $L$).
\item $O$ decomposes $w=xyz$ with $|xy|\le m$, $|y|\ge1$.
\item $P$ picks $i$ such that $xy^iz\notin L$. If such $i$ exists for \emph{every} decomposition, $L$ is not regular.
\end{enumerate}
\ex{4.7} $L=\{a^nb^n:n\ge0\}$ not regular. Given $m$, pick $w=a^mb^m$. Since $|xy|\le m$, $y=a^k,\ k\ge1$. Pump $i=2$: $xy^2z=a^{m+k}b^m\notin L$. Contradiction.
\ex{4.8} $L=\{ww^R:w\in\Sigma^*\}$ not regular. Pick $w=a^mb^mb^ma^m$ (in $L$); $y=a^k$; $i=2$: $a^{m+k}b^{2m}a^m\notin L$ (not a palindrome form). \emph{Pitfall:} picking $a^{2m}$ lets opponent win ($y=aa$ pumps fine) -- choose $w$ that kills every allowed $y$.
\ex{4.9} $L=\{w:n_a(w)<n_b(w)\}$ not regular: $w=a^mb^{m+1}$, $y=a^k$, $i=2$: $a^{m+k}b^{m+1}\notin L$.
\ex{4.10} $L=\{(ab)^na^k:n>k\ge0\}$ not regular: $w=(ab)^{m+1}a^m$; $y$ inside first $m$ chars; $i=0$: $(ab)^pb(ab)^qa^m$ or $(ab)^ma^m\notin L$.
\ex{4.11} $L=\{a^{n^2}:n\ge0\}$ not regular: $w=a^{m^2}$, $y=a^k,1\le k\le m$; $i=0$: $a^{m^2-k}\notin L$ since $(m-1)^2<m^2-k<m^2$.
\ex{4.12} $L=\{a^nb^kc^{n+k}:n,k\ge0\}$ not regular: $w=a^mb^mc^{2m}$, $i=0$: $a^{m-k}b^mc^{2m}\notin L$. (Or: homomorphism $h(a)=a,h(b)=a,h(c)=c$ maps $L$ to $\{a^{n+k}c^{n+k}\}$.)
\ex{4.13} $L=\{a^nb^l:n\neq l\}$ not regular: $w=a^{m!}b^{(m+1)!}$, $y=a^k,1\le k\le m$; need $i$ with $m!+(i-1)k=(m+1)!$: $i=1+\frac{m\cdot m!}{k}$, integer since $k\le m$. (Also via $L\cap a^*b^*$.)
\ex{} $L=\{a^p:p$ prime$\}$ not regular: $w=a^p$, $p$ = smallest prime $\ge m$; $y=a^k$; choose $i-1=p$: $xy^{p+1}z=a^{p+pk}=a^{p(k+1)}$ composite.
\textbf{Common pitfalls:} (1) Pumping lemma can \emph{never} prove regularity (it is one-way; e.g.\ $\{a^ib^jc^k: i=1\Rightarrow j=k\}$ satisfies PL but is nonregular); (2) starting with $w\notin L$ (e.g.\ $a^m$ for the primes language); (3) assuming a specific decomposition instead of handling \emph{all} decompositions allowed by $|xy|\le m$.

%======================================================================
\section{Ch 5 -- Context-Free Grammars \& Parsing (Lec 5)}
%----------------------------------------------------------------------
\subsection{Context-Free Grammars}
\df{5.1 CFG} $G=(V,T,S,P)$; productions $A\to x$ with $A\in V$, $x\in(V\cup T)^*$: the substitution of $A$ can happen \emph{regardless of context}.
\df{} $L(G)=\{w\in T^*: S\overset{*}{\Rightarrow}w\}$; $L$ is a CFL iff some CFG generates it.
Regular $\subsetneq$ CFL: e.g.\ $\{a^nb^n\}$, $\{ww^R\}$ are CFL but not regular.
\ex{5.1} $S\to aSa|bSb|\lambda$: $L=\{ww^R:w\in\{a,b\}^*\}$ (even palindromes).
\ex{5.2} $S\to abB,\ A\to aaBb,\ B\to bbAa,\ A\to\lambda$: $L=\{ab(bbaa)^nbba(ba)^n:n\ge0\}$.
\ex{5.3} $\{a^nb^m:n\neq m\}$ is CFL: $S\to AS_1|S_1B$ (n>m / n<m), $S_1\to aS_1b|\lambda$, $A\to aA|a$, $B\to bB|b$. (Not linear.)
\ex{5.4} $S\to aSb|SS|\lambda$: $L=\{w:n_a(w)=n_b(w)$ and every prefix has $n_a\ge n_b\}$ = matched parentheses. (Note: adding $S\to bSa$ gives the unconstrained equal-count language.)
\textbf{Derivation order:} leftmost = always replace leftmost variable; rightmost = always replace rightmost variable. In linear grammars the order is forced; in general CFGs it is not.
\ex{5.5} $S\to aAB,\ A\to bBb,\ B\to A|\lambda$. Leftmost: $S\Rightarrow aAB\Rightarrow abBbB\Rightarrow abAbB\Rightarrow abbBbbB\Rightarrow abbbbB\Rightarrow abbbb$. Rightmost: $S\Rightarrow aAB\Rightarrow aA\Rightarrow abBb\Rightarrow abAb\Rightarrow abbBbb\Rightarrow abbbb$.
\df{5.3 derivation (parse) tree} ordered tree: (1) root $S$; (2) leaves from $T\cup\{\lambda\}$; (3) internal vertices from $V$; (4) node $A$ with children $a_1\cdots a_n$ requires production $A\to a_1\cdots a_n$; (5) a $\lambda$-leaf has no sibling. \textbf{Yield} = leaves read left-to-right. Partial derivation tree: properties 3--5, root may be any variable, leaves in $V\cup T\cup\{\lambda\}$; yield is a sentential form.
\thm{5.1} $w\in L(G)$ iff there is a derivation tree with yield $w$; yield of any partial tree rooted at $S$ is a sentential form.
\textbf{One derivation tree $\leftrightarrow$ many derivations:} leftmost and rightmost derivations of $aab$ from $S\to AB, A\to aaA|\lambda, B\to Bb|\lambda$ give the same tree.

%----------------------------------------------------------------------
\subsection{Parsing}
Parsing = given $w$, find a derivation sequence of productions.
\textbf{Exhaustive search (top-down):} round 1: all $S\Rightarrow x$ one-step forms; round $k$: apply all productions to each form of round $k{-}1$ (e.g.\ $|P|^k$ forms). Flaws: tedious; may \emph{not terminate} for $w\notin L(G)$ when $\lambda$ or unit productions exist (sentential form length doesn't grow).
\ex{5.7} $S\to SS|aSb|bSa|\lambda$, find $aabb$: phase 3 finds $S\Rightarrow aSb\Rightarrow aaSbb\Rightarrow aabb$.
\ex{5.8} Replace by $S\to SS|aSb|bSa|ab|ba$ (no $\lambda$, no unit): each round lengthens the form, so parsing $w$ ends within $|w|$ rounds.
\thm{5.2} If CFG has no $A\to\lambda$ and no $A\to B$ productions, exhaustive search terminates for \emph{any} $w$: $\le 2|w|$ phases, $\le |P|+|P|^2+\cdots+|P|^{2|w|}$ forms $\Rightarrow O(|P|^{2|w|+1})$ -- terrible. General CFG membership: $O(|w|^3)$ (CYK, Lec 6).
\textbf{S-grammar (simple grammar):} productions $A\to ax$ ($a$ terminal, $x\in V^*$) and each pair $(A,a)$ occurs \emph{once}. Parsing is deterministic: one choice per phase, total time $|w|$ (linear).
\ex{5.9} $S\to aS|bSS|c$ is an s-grammar: $S\Rightarrow aS\Rightarrow abSS\Rightarrow abcS\Rightarrow abcc$ unique.

%----------------------------------------------------------------------
\subsection{Ambiguity}
\df{5.5} CFG $G$ is \textbf{ambiguous} iff some $w\in L(G)$ has $\ge 2$ distinct derivation trees, equivalently $\ge 2$ leftmost (or rightmost) derivations.
\ex{5.11} $E\to E+E|E*E|(E)|a$: $a+a*a$ has two trees ($(a+a)*a$ vs $a+(a*a)$) $\Rightarrow$ ambiguous. With values $2+2*2$: results 8 vs 6 -- ambiguity is dangerous for PLs.
\textbf{Fix (unambiguous precedence grammar):}
$E\to E+T|T,\quad T\to T*F|F,\quad F\to(E)|a$. $*$ binds tighter than $+$; unique tree for $a+a*a$ giving $a+(a*a)$.
\textbf{Dangling else:} IF\_STMT $\to$ if E then S | if E then S else S is ambiguous (two trees for ``if e1 then if e2 then s1 else s2''). Fix by matching else with nearest unmatched then.
\textbf{Inherent ambiguity:} some CFLs have \emph{only} ambiguous grammars, e.g.\ $L=\{a^nb^nc^m\}\cup\{a^nb^mc^m\}$: string $a^nb^nc^n$ always has two trees. (Deciding inherent ambiguity: hard.)
\ex{Ex 5.3.6} $S\to AB|aaB,\ A\to a|Aa,\ B\to b$ ambiguous: two leftmost derivations of $aab$: $S\Rightarrow aaB\Rightarrow aab$ and $S\Rightarrow AB\Rightarrow AaB\Rightarrow aaB\Rightarrow aab$.
\textbf{CFGs and programming languages:} BNF (Backus-Naur form) $\langle expr\rangle::=\langle term\rangle|\langle expr\rangle+\langle term\rangle$; LL and LR grammars parse in linear time. Syntactic (grammar) vs semantic (types) errors are separate.

%======================================================================
\section{Ch 6 -- Grammar Transformations \& Normal Forms (Lec 6)}
%----------------------------------------------------------------------
\subsection{Transformation Methods}
\textbf{Substitution rule} \thm{6.1}: if $A\to x_1Bx_2$ and $B\to y_1|\cdots|y_n$ (all $B$-productions), replace by $A\to x_1y_1x_2|\cdots|x_1y_nx_2$ (optionally keep $A\to x_1Bx_2$). $L$ unchanged.
\ex{6.1} $A\to a|aaA|abBc,\ B\to abbA|b$ $\Rightarrow$ $A\to a|aaA|ababbAc|abbc$.
\textbf{(1) Remove $\lambda$-productions} ($A\to\lambda$): find all \emph{nullable} variables ($A\overset{*}{\Rightarrow}\lambda$); for each production, add all variants with each nullable variable kept or removed (skip the variant that deletes everything if it gives $A\to\lambda$ for non-start). Keep $S\to\lambda$ only if $\lambda\in L(G)$ (via extra start $S_0\to S|\lambda$).
\ex{} $S\to aMb,\ M\to aMb|\lambda$: nullable $\{M\}$; result $S\to aMb|ab,\ M\to aMb|ab$.
\ex{6.5} $S\to ABaC,\ A\to BC,\ B\to b|\lambda,\ C\to D|\lambda,\ D\to d$; nullable $\{A,B,C\}$: $S\to ABaC|BaC|AaC|ABa|aC|Aa|Ba|a$; $A\to B|C|BC$; $B\to b$; $C\to D$; $D\to d$.
\textbf{(2) Remove unit-productions} ($A\to B$): remove $A\to A$ immediately; for each $B\to y_1|\cdots|y_n$ (non-unit), add $A\to y_1|\cdots|y_n$; iterate (use dependency graph). 
\ex{6.6} $S\to Aa|B,\ B\to A|bb,\ A\to a|bc|B$: dependency $A\to B\to A$; final $S\to a|bc|bb|Aa$, $A\to a|bb|bc$, $B\to a|bb|bc$.
\textbf{(3) Remove useless productions:} variable useful iff $S\overset{*}{\Rightarrow}xAy\overset{*}{\Rightarrow}w\in T^*$. Two passes: (a) find variables that derive terminal strings (bottom-up marking), delete others \emph{and their productions}; (b) draw dependency graph (edge $A\to B$ iff production $A\to xBy$); delete variables unreachable from $S$. Production useless iff it contains a useless variable.
\ex{6.3} $S\to aS|A|C,\ A\to a,\ B\to aa,\ C\to aCb$: pass (a): $\{A,B\}$ derive terminals; $S\to A$ adds $S$; $C$ never terminates $\Rightarrow$ delete $C$: $S\to aS|A,\ A\to a,\ B\to aa$. Pass (b): $B$ unreachable $\Rightarrow$ delete: $S\to aS|A,\ A\to a$.
\thm{6.5} If $L$ is CFL without $\lambda$: exists CFG without useless/unit/$\lambda$-productions. \textbf{Removal order matters:} (1) $\lambda$-productions, (2) unit-productions, (3) useless productions.

%----------------------------------------------------------------------
\subsection{Normal Forms}
\df{CNF (Chomsky)}: all productions $A\to BC$ (two variables) or $A\to a$ (one terminal). (May need $S\to\lambda$ handled separately.)
\textbf{Why CNF:} parse trees become binary; for a string of length $n$, tree depth $\ge\log n$; enables pumping lemma for CFLs and CYK.
\ex{6.7} $S\to AS|a,\ A\to SA|b$ is CNF. ($S\to AAS$, $A\to aa$ are not.)
\textbf{CNF procedure} \thm{6.6}:
\begin{enumerate}
\item Remove $\lambda$-productions, unit-productions, useless productions.
\item For each terminal $a$: add $T_a\to a$ and replace $a$ by $T_a$ in all productions with RHS length $\ge2$.
\item Replace $A\to C_1C_2\cdots C_n$ ($n\ge3$) by $A\to C_1V_1,\ V_1\to C_2V_2,\ \dots,\ V_{n-2}\to C_{n-1}C_n$.
\end{enumerate}
\ex{6.8} $S\to ABa,\ A\to aab,\ B\to Ac$: terminals: $T_a\to a,T_b\to b,T_c\to c$; $S\to ABT_a,\ A\to T_aT_aT_b,\ B\to AT_c$; binarize: $S\to AV_1,\ V_1\to BT_a$; $A\to T_aV_2,\ V_2\to T_aT_b$; $B\to AT_c$.
\df{GNF (Greibach)}: all productions $A\to aV_1\cdots V_k$, $k\ge0$ (terminal first, then only variables).
\ex{} $S\to cAB,\ A\to aA|bB|b,\ B\to b$ is GNF. $S\to abSb$, $S\to aa$ is not.
\ex{6.9} $S\to AB,\ A\to aA|bB|b,\ B\to b$ $\Rightarrow$ $S\to aAB|bBB|bB$.
\ex{6.10} $S\to abSb|aa$ $\Rightarrow$ $S\to aT_bST_b|aT_a,\ T_a\to a,\ T_b\to b$.
\thm{6.7} Every CFG without $\lambda$ has an equivalent GNF. \textbf{Observation:} GNF is excellent for parsing but hard to compute; CNF is easy to find.

%----------------------------------------------------------------------
\subsection{CYK Membership Algorithm}
Input: CFG in CNF, string $w$. Output: is $w\in L(G)$?
\textbf{Dynamic programming over substrings:} let $V_{ij}=\{A: A\overset{*}{\Rightarrow}w_i\cdots w_j\}$ (variables deriving substring $w_i\cdots w_j$).
\begin{itemize}
\item \textbf{Base} (length 1): $V_{ii}=\{A: A\to w_i\in P\}$.
\item \textbf{Induction} (length $>1$): for all $k$, $i\le k<j$: if $A\to BC\in P$, $B\in V_{ik}$, $C\in V_{k+1,j}$, then add $A$ to $V_{ij}$.
\item $w\in L(G)$ iff $S\in V_{1n}$.
\end{itemize}
\textbf{Time:} $O(|w|^3)$ (positions $O(n^2)$, splits $O(n)$); also converts to a bottom-up parser.
\ex{} CNF $S\to AB,\ A\to BB|a,\ B\to AB|b$; $w=aabbb$:
length-1: $V_{11}=V_{22}=\{A\}$, $V_{33}=V_{44}=V_{55}=\{B\}$.
length-2: $aa\mapsto\emptyset$; $ab\mapsto\{S,B\}$ (both $S\to AB$ and $B\to AB$ with $A\in V_{11},B\in V_{33}$); $bb\mapsto\{A\}$ ($A\to BB$), so $V_{34}=V_{45}=\{A\}$.
length-3: $aab\mapsto\{S,B\}$; $abb\mapsto\{A\}$; $bbb\mapsto\{S,B\}$.
length-4: $aabb\mapsto\{A\}$; $abbb\mapsto\{S,B\}$.
length-5: $aabbb\mapsto\{S,B\}\ni S$ $\Rightarrow$ $aabbb\in L(G)$.
\textbf{Layout as a triangle} above the string (as in slides) with cells $V_{ij}$.

%======================================================================
\section{Quick Reference: Theorems \& Definitions}
\fbox{\parbox{\dimexpr\linewidth-2\fboxsep-4pt}{\scriptsize
\begin{tabular}{@{}p{0.29\linewidth}p{0.63\linewidth}@{}}
\setlength{\tabcolsep}{3pt}
Def 2.1 & DFA $M=(Q,\Sigma,\delta,q_0,F)$, total $\delta$\\
Def 2.5/2.6 & NFA: $\delta:Q\times(\Sigma\cup\{\lambda\})\to2^Q$; $L(M)=\{w:\delta^*(q_0,w)\cap F\ne\emptyset\}$\\
Thm 2.2 & NFA languages = regular languages (subset construction)\\
Def 2.8 & indistinguishable states $\Rightarrow$ DFA minimization\\
Def 3.1 & REs; primitive $\emptyset,\lambda,a$; $+,\cdot,^*$\\
Thm 3.1/3.2 & RE $\Leftrightarrow$ regular (Kleene): NFA constructions / GTG elimination\\
Thm 3.3/3.4 & Regular grammars $\Leftrightarrow$ regular languages\\
Thm 4.8 & Pumping Lemma: $w=xyz$, $|xy|\le m$, $|y|\ge1$, $xy^iz\in L\ \forall i$\\
Thm 4.3 & Homomorphic image of regular language is regular\\
Def 5.1 & CFG $G=(V,T,S,P)$, $A\to x$, $x\in(V\cup T)^*$\\
Def 5.5 & ambiguity: $\ge2$ trees / leftmost derivations\\
Thm 6.1 & substitution rule preserves $L(G)$\\
Thm 6.5 & remove $\lambda$, then unit, then useless productions\\
Thm 6.6 & every CFG (no $\lambda$) $\to$ CNF\\
Thm 6.7 & every CFG (no $\lambda$) $\to$ GNF\\
CYK & membership in CNF, dynamic programming, $O(n^3)$\\
\end{tabular}}}
\end{multicols*}
\end{document}
